%% Section 2: Model and Problem Formulation. Owns all math objects.
%% Notation locked in outline v2 Section 1 ledger.

\begin{figure}[t]
  \centering
  \figplaceholder{F1 diptych (outline v2 sec.4): (a) ordinary batching ---
    audio revision, locally differing batched proposal displayed, output
    enters feedback state, trajectories diverge; corner annotation
    ``8/21 legal pairs propagate with feedback; 0/21 without; tail RAF up
    to 153 measured on the 6 persistent pairs''. (b) DGFT on the identical
    scenario --- local difference retained at display, state write denied,
    coarse authority later produces certificate, frontier advances.}
  {State-entry contamination under ordinary batching (a) and the same
  scenario under DGFT (b). DGFT does not require every provisional output
  to be identical; it prevents uncertified outputs from becoming
  future-state parents.}
  \label{fig:motivation}
\end{figure}

\subsection{Revision-generative state machines}\label{sec:model:rsm}

\begin{definition}[Revision-generative serving round]\label{def:rsm}
A serving round $t$ of a revision-generative system is a tuple
$(\state_t,\audio_t,\out_t)$ where $\state_t$ is the committed feedback
state entering the round, $\audio_t$ is the exogenous input (audio frontier),
and $\out_t$ is the generated output. The state update is
\begin{equation}\label{eq:state-update}
  \state_{t+1} = \mathcal{T}(\state_t,\audio_t,\out_t),
\end{equation}
where $\mathcal{T}$ is the model/adapter state-update function. The system is
\emph{revision-generative} when $\mathcal{T}$ depends on $\out_t$, i.e.\ the
output of round $t$ participates in the state read by round $t{+}1$.
\end{definition}

Not every byte of $\out_t$ enters future state. We isolate exactly the part
that does:

\begin{definition}[State projection]\label{def:projection}
The state projection $\proj(\out_t;\state_t)$ is the component of $\out_t$
that $\mathcal{T}$ consumes when forming $\state_{t+1}$. For Qwen3-ASR it is
the rollback prefix re-prompted at the next revision~\cite{shi2026qwen3asr};
for Whisper with previous-text conditioning it is the conditioning
prompt~\cite{radford2023whisper}.
\end{definition}

Ordinary batched serving makes an implicit assumption we name explicitly:
\begin{equation}\label{eq:implicit}
  \state_{t+1} = \proj(\out_t^{\mathrm{batch}};\state_t),
\end{equation}
i.e.\ whatever the batched forward pass produced may write state.
Equation~\eqref{eq:implicit} conflates ``was generated'' with ``may write
future state''.

\subsection{Display output versus state projection}\label{sec:model:display}

The user-facing display function may surface output that carries no
future-state authority: display-only suffixes, hold-back regions, and
provisional hypotheses are shown but excluded from
$\proj(\out_t;\state_t)$ by the adapter (Section~\ref{sec:system}). The
paper's central invariant separates the two permissions:
\begin{equation}\label{eq:invariant}
  \boxed{\;\mathrm{Display}(\out_t)\;\not\Rightarrow\;
  \mathrm{Authorize}\bigl(\proj(\out_t;\state_t)\bigr)\;}
\end{equation}

\subsection{State-entry divergence and contamination}\label{sec:model:divergence}

Let $\out_t^{f}$ be the fine (batched, proposer) output and $\out_t^{a}$ the
canonical authority output for the same round.

\begin{definition}[Local output divergence]\label{def:eps}
$\varepsilon_t = d_{\out}(\out_t^{f},\out_t^{a})$, with $d_{\out}$ an
output-space distance (mixed-unit edit distance in our contracts).
\end{definition}

\begin{definition}[State-entry divergence]\label{def:e}
$\ediv_t = d_{\state}\bigl(\proj(\out_t^{f};\state_t),\,
\proj(\out_t^{a};\state_t)\bigr)$, with $d_{\state}$ the state-space
distance induced by the same mixed-unit metric.
\end{definition}

\begin{definition}[Contamination event]\label{def:contamination}
Round $t$ is a contamination event $\cont_t$ iff the state projection of a
fine proposal is installed into committed state while $\ediv_t > 0$:
\[
  \cont_t \;:\Longleftrightarrow\; \ediv_t > 0 \;\wedge\;
  \proj(\out_t^{f};\state_t) \text{ installed into } \state_{t+1}.
\]
\end{definition}

Definition~\ref{def:contamination} upgrades the problem statement from
``outputs differ'' to ``a differing output acquired future-state write
permission''. Display-level instability literature observes the former;
this paper targets the latter.

\subsection{Feedback amplification}\label{sec:model:amplification}

Let $\dist_t = d_{\state}(\state_t^{f},\state_t^{a})$ be the trajectory
distance between the batched-served state trajectory and the canonical one.
A local perturbation decomposition gives
\begin{equation}\label{eq:recurrence}
  \dist_{t+1} \;\le\; L_t\,\dist_t \;+\; \beta_t\,\ediv_t ,
\end{equation}
where $L_t$ measures propagation of pre-existing difference and $\beta_t$
measures coupling of the round-$t$ state-entry divergence into the next
state. We make \emph{no} global Lipschitz claim: \eqref{eq:recurrence} is a
per-round decomposition with unknown coefficients, used only to identify
channels.

\begin{remark}[Causal prediction]\label{rem:causal}
Removing the feedback edge sets $\beta_t \equiv 0$, so a local state-entry
divergence can no longer enter $\dist_{t+1}$ through the revision channel.
Section~\ref{sec:eval} (H1) tests this prediction by intervention.
\end{remark}

\begin{definition}[Revision amplification factor]\label{def:raf}
For horizon $K$, the measured trajectory amplification factor is
\[
  \mathrm{RAF}_{t,K} \;=\; \frac{\dist_{t+K}}{\ediv_t + \epsilon_0},
\]
an empirical indicator of how one state-entry divergence propagates through
subsequent revisions. RAF is reported as a measured trajectory quantity, not
as a theoretical bound.
\end{definition}

\subsection{Design requirements derived from the model}\label{sec:model:reqs}

The model above forces four requirements on any serving mechanism that
claims to certify feedback state:
\begin{description}
  \item[R1 (Authority separation).] Display permission and state-write
    permission are distinct grants; Eq.~\eqref{eq:invariant}.
  \item[R2 (Parent ownership).] A state projection commits only onto the
    parent version it was computed against; otherwise staleness is
    undetectable.
  \item[R3 (Frontier monotonicity).] Committed audio frontiers advance
    monotonically---no gap, no overlap---so the committed history is a
    single sequence.
  \item[R4 (Evidence non-authority).] Statistical evidence (margins,
    coverage, repetition signals) schedules verification; it never
    authorizes installation.
\end{description}
R1--R4 reappear as protocol properties of DGFT
(Section~\ref{sec:dgft}), labels on the system diagram
(Figure~\ref{fig:arch}), and columns of the guarantee taxonomy
(Table~\ref{tab:guarantees}).
